\section{Arbitrary termination in dimension three}\label{sec:threefold}

We establish the lower-dimensional termination needed for special
termination.  Throughout this section we use the rational generalized
pairs and elementary steps of Section~\ref{sec:relative}.
The spaces are compact K\"ahler and globally
strongly $\Q$-factorial, and the elementary steps have the negative-side
one-ray condition for degrees of global rational line bundles.  All higher
nef data are fixed as a rational b-line bundle.  The projective bases of the
steps may change.  No pseudo-effectivity assumption is made in this section.
We first treat generalized terminal pairs.  We then bound Cartier indices
under uniform discrepancy bounds and use that bound to lift a general
sequence to generalized terminal models.  This follows the threefold
strategy of \cite[Section~3]{CT2023}; the surface, index, and lifting
arguments below establish its analytic form with varying bases.

\subsection{Generalized terminal pairs}

\begin{proposition}\label{prop:terminal-three}
Let $(X_0,B_0+\bM)$ be a rational generalized terminal pair on a normal
globally strongly $\Q$-factorial compact K\"ahler threefold, with fixed
rational b-line data as in Definition~\ref{def:generalized-data}.
Every sequence of small elementary steps in the sense of
Section~\ref{sec:relative} starting from
this pair terminates.  The contraction bases may vary.  Here generalized
terminal means generalized klt with all exceptional log discrepancies
greater than one.
\end{proposition}

\begin{proof}
Write the sequence as
\[
 (X_i,B_i+\bM)\dashrightarrow (X_{i+1},B_{i+1}+\bM),
 \qquad X_i\xrightarrow{f_i}Z_i\xleftarrow{f_i^+}X_{i+1}.
\]
The boundary is transformed strictly, so its coefficients belong to one
finite set.  Smallness and Lemma~\ref{lem:negativity} preserve generalized
terminality.  Forgetting the effective boundary and the effective nef defect
increases discrepancies, so each $X_i$ is an ordinary terminal threefold.
In particular its singularities are isolated, and the general point of every
flipped curve is smooth.

Choose an integer $p>0$ for which $p\bM$ is represented by an integral line
bundle on a carrier, and an integer $N$ divisible by $p$ and by the
denominators of the boundary coefficients.  These integers continue to work
on higher carriers throughout any finite part of the sequence.  If $C$ is
a flipped curve on $X_{i+1}$, the exceptional valuation $E_C$ of its generic
blowup has log discrepancy
\begin{equation}\label{eq:curve-discrepancy-three}
 a(E_C;X_{i+1},B_{i+1}+\bM)
   =2-\sum_j m_jb_j-\frac{q}{p},
 \qquad m_j,q\in\Z_{\geq0}.
\end{equation}
Indeed, the ordinary blowup contribution is two, the $m_j$ are the generic
multiplicities of the boundary components, and the remaining term is the
coefficient of the nef defect.  At this smooth generic point the integral
trace of $p\bM$ is Cartier.  The defect therefore has coefficients in
$p^{-1}\Z$ and is effective by Lemma~\ref{lem:trace}.  This proves
\eqref{eq:curve-discrepancy-three}, including its signs.

Let $b$ be the largest boundary coefficient, with $b=0$ when the boundary
is empty.  For $b>0$ test the flipped curves contained in a coefficient-$b$
component; for $b=0$ test every flipped curve.  The value in
\eqref{eq:curve-discrepancy-three} belongs to the fixed finite set
\[
 \mathcal A_b=N^{-1}\Z\cap(0,2-b].
\]
For each $t\in\mathcal A_b$, set
\[
 d_t(i)=\#\{E\text{ exceptional over }X_i:
                   a(E;X_i,B_i+\bM)<t\}.
\]
Lemma~\ref{lem:low-places} makes these numbers finite, since the terminal
cutoff is the strict inequality $a<2-b$.  Small transformations identify
the sets of exceptional valuations, and discrepancy monotonicity makes
every $d_t(i)$ nonincreasing.  For the valuation of a tested curve,
Lemma~\ref{lem:negativity} gives
\[
 a(E_C;X_i,B_i+\bM)
 <a(E_C;X_{i+1},B_{i+1}+\bM)=t.
\]
Thus $d_t$ drops at that step.  The nonnegative integer
$\sum_{t\in\mathcal A_b}d_t(i)$ can drop only finitely often.  If $b=0$,
every nontrivial small step has a flipped curve, so the proof is complete.

Suppose $b>0$.  On a tail, no flipped curve is contained in any of the
coefficient-$b$ components.  Fix such a component $D_i$, and let
$D_i^\nu$ denote its normalization.  The two images of $D_i$ and
$D_{i+1}$ in $Z_i$ agree, because the ambient transformation is an
isomorphism in codimension one.  Denote the normalization of this common
image by $V_i$.  There are projective birational morphisms
\[
 D_i^\nu\longrightarrow V_i\longleftarrow D_{i+1}^\nu.
\]
The morphism on the right contracts no curve and is therefore finite:
every positive-dimensional projective fiber contains a curve.  A finite
birational morphism to a normal space is an isomorphism.  Hence there is a
projective birational morphism $D_i^\nu\to D_{i+1}^\nu$.  It contracts
exactly the curves whose images in $X_i$ are flipping curves contained in
$D_i$.  By Lemma~\ref{lem:cycle-loss}, with $k=1$, each such contraction
strictly decreases the dimension of the span of curve classes of the compact
normalized surface.  There are finitely many coefficient-$b$ components,
so only finitely many further steps have a flipping curve contained in one
of them.

On the resulting tail put $D_i^{\max}=\sum_{b_j=b}D_{i,j}$ and replace
$B_i$ by $B_i-bD_i^{\max}$.  Every contracted curve $C$ has
$D_i^{\max}\cdot C\geq0$, because it is not contained in this effective
$\Q$-Cartier divisor.  Thus the new adjoint
\[
 A_i'=K_{X_i}+B_i-bD_i^{\max}+M_{X_i}
\]
is negative on the old ray.  To check its positive transform, write
$A_i=K_{X_i}+B_i+M_{X_i}$ and choose the rational number
$c=(A_i'\cdot C)/(A_i\cdot C)>0$.  The global one-ray condition and
Lemma~\ref{lem:descent} give an actual identity of rational line bundles
\[
 A_i'=cA_i+f_i^*L_i,\qquad L_i\in\Pic(Z_i)\otimes\Q.
\]
Its strict transform is
$A_{i+1}'=cA_{i+1}+(f_i^+)^*L_i$, which is $f_i^+$-ample.
Consequently the same tail is a sequence of elementary steps for the new
generalized pair.  Decreasing the effective boundary preserves generalized
terminality.  Induction on the number of distinct positive boundary
coefficients proves the proposition.
\end{proof}

\subsection{Uniform indices for the terminal lift}

To lift a generalized klt sequence to generalized terminal models, we
will need the discrepancies of the finitely many exceptional places with
log discrepancy at most one to stabilize.  A uniform Cartier index will
put those discrepancies in a fixed rational grid.  We first establish
the local index statement and the surface estimate used to obtain that
uniform bound.

\begin{lemma}\label{lem:terminal-local-index}
Let $(T,x)$ be a complex analytic terminal threefold germ of canonical
Cartier index $r$.  If $D$ is an integral Weil divisor germ that is
$\Q$-Cartier, then $rD$ is Cartier.  If $r>1$, there is an exceptional
valuation centered at $x$ with ordinary log discrepancy $1+1/r$.
\end{lemma}

\begin{proof}
For the first assertion, use the analytic index-one cover
$\pi:(T^\sharp,x^\sharp)\to(T,x)$, a cyclic cover of degree $r$ that is
\'etale away from the distinguished point.  The terminal threefold
classification makes $T^\sharp$ smooth or an isolated cDV hypersurface
germ; see \cite{Mori1985}.  A sufficiently small punctured representative
$U^\sharp$ is simply connected, by the connectivity theorem for links of
isolated hypersurfaces \cite[Theorems~2.10 and~5.2]{Milnor1968}.  For a convergent
analytic defining equation, finite determinacy gives an analytically
equivalent sufficiently high Taylor polynomial
\cite[Theorem~6.2, p.~24]{DeJongEtAl1998}, so the same connectivity
statement applies; the smooth case has the same property.  For every
$m>0$, the analytic Kummer sequence implies
that the $m$-torsion of $\Pic(U^\sharp)$ is an image of
$H^1(U^\sharp,\boldsymbol\mu_m)=0$.  The local divisor class group
injects into this Picard group: a rank-one reflexive sheaf trivial on the
punctured normal germ is trivial on the germ by reflexive extension.
It follows that the local class group of $T^\sharp$ is torsion-free.

Since $D$ is $\Q$-Cartier, the integral divisor $\pi^*D$ has torsion class
and is therefore principal, say $\pi^*D=\operatorname{div}(u)$ for a
meromorphic germ $u$.  The meromorphic norm gives
\[
 \operatorname{div}\operatorname{Norm}_{\pi}(u)
   =\pi_*\operatorname{div}(u)
   =\pi_*\pi^*D=rD.
\]
Thus $rD$ is Cartier; compare the index divisibility statement in
\cite[Lemma~5.1]{Kawamata1988}.  The argument above proves the analytic
assertion for divisor germs already known to be $\Q$-Cartier.

The second assertion is the terminal small-discrepancy theorem
\cite[Appendix, pp.~193--195]{KawamataAppendix1992}.  Its proof uses the
analytic quotient classification and gives ordinary discrepancy $1/r$, or
log discrepancy $1+1/r$ in our convention.
\end{proof}

When the germ lies on a compact terminal space, the last valuation can be
seen on a global resolution.  In fact, its ordinary crepant boundary has
only negative exceptional coefficients.  A valuation still exceptional
over that smooth resolution has ordinary log discrepancy at least two.
Thus a local valuation of discrepancy $1+1/r<2$ is already represented by
an exceptional component of the restricted global resolution.  Since its
center on $T$ is the isolated point $x$, that component gives a global
exceptional prime.  This observation permits the use of
Lemma~\ref{lem:terminal-local-index} with global discrepancy bounds.

\begin{lemma}\label{lem:exceptional-surface-curve}
Let $h:U\to T$ be a projective bimeromorphic morphism of normal compact
complex threefolds.  Suppose that $U$ is klt and that $P$ is an exceptional
prime divisor with $K_U+P$ rational Cartier.  Given finitely many effective
$\Q$-Cartier divisors $D_j$ on $U$, none having $P$ as a component, there
is a curve $C\subset P$ contracted by $h$, contained in none of the $D_j$,
such that
\[
 (K_U+P)\cdot C\geq-3.
\]
\end{lemma}

\begin{proof}
Let $\nu:P^\nu\to P$ be the normalization and let
$\mu:S\to P^\nu$ be its minimal resolution.  Divisorial adjunction gives
\begin{equation}\label{eq:arbitrary-prime-different}
 \nu^*((K_U+P)|_P)=K_{P^\nu}+\Delta_{P^\nu},
 \qquad \Delta_{P^\nu}\geq0,
\end{equation}
as an identity of rational adjoint bundles.  Neither normality of $P$ nor
log canonicity of $(U,P)$ is needed for effectivity of the different in this
formula.  This is the ordinary different of a reduced prime, computed at
codimension-two points of $U$; the construction and effectivity are given in
\cite[Section~14.1(i)--(iii), pp.~40--41]{FujinoFundamental2010},
and analytic adjunction is included in
\cite[Theorem~1.1 and the following setup, p.~1]{FujinoInversion2023}.
Here is the local analytic effectivity argument.  On a transverse klt
surface $T$, take a small smooth quotient $q:V\to T$ and write
$P_T$ for the reduced boundary curve and $C=(q^{-1}P_T)_{\mathrm{red}}$.
The quotient is unramified at the general points of these curves, so
$q^*(K_T+P_T)=K_V+C$.  Let $\rho:C^\nu\to P_T^\nu$ be the induced
map of normalizations.  Residue adjunction for the reduced plane curve
$C$ has an effective conductor divisor $A$, and the canonical
ramification formula for $\rho$ has an effective divisor $R$.
Compatibility of residues with pullback gives
\[
 \rho^*\operatorname{Diff}_{P_T^\nu}(0)=A+R.
\]
At a branch of ramification degree $e$, its coefficient downstairs is
$(c+e-1)/e\geq0$, where $c$ is the conductor coefficient.  This proves
the analytic effectivity in \eqref{eq:arbitrary-prime-different}
without an lc hypothesis on $(U,P)$ or normality of $P$.

The sum on the right of \eqref{eq:arbitrary-prime-different} is rational
Cartier; its individual canonical term need not be.  If that term is used
separately, take its numerical pullback in the sense of Mumford, determined
by intersection zero with all exceptional curves.  The negative definite
exceptional intersection matrix makes this pullback linear and makes the
pullback of an effective divisor effective.  In particular the calculation
does not presume $\Q$-factoriality of $P^\nu$.

For completeness, the effective correction can also be checked directly
on $S$.  For every irreducible $\mu$-exceptional curve $E$, adjunction on
the smooth surface gives
\[
 K_S\cdot E=2p_a(E)-2-E^2\geq0.
\]
Indeed, an exceptional smooth rational $(-1)$-curve is excluded by
minimality over $P^\nu$; every other exceptional curve satisfies the
displayed inequality.  The pullback of the adjoint in
\eqref{eq:arbitrary-prime-different} therefore has the form
\begin{equation}\label{eq:surface-effective-correction}
 (\nu\mu)^*((K_U+P)|_P)=K_S+\Delta_S,
 \qquad \Delta_S\geq0.
\end{equation}
To obtain the sign, the residue comparison defines $\Delta_S$ as a rational
divisor with pushforward $\Delta_{P^\nu}\geq0$.  Since $K_S$ is
$\mu$-nef, the divisor $-\Delta_S$ is $\mu$-nef.  Apply
Lemma~\ref{lem:negativity}\textup{(i)} to $E=\Delta_S$ to obtain
$\Delta_S\geq0$.  This is equivalently the
numerical-pullback calculation above, and proves
\eqref{eq:surface-effective-correction} as an actual bundle comparison.

The image $h(P)$ has dimension zero or one.  In the latter case, general
smooth fiber components of $S\to h(P)$ are moving curves $F$ with
$K_S\cdot F=2g(F)-2\geq-2$.  One may use the Stein factorization to
describe these fibers over a smooth compact curve.  Choose a general fiber
component not contained in $\Supp\Delta_S$, the conductor, or the
pullbacks of the $D_j$.

If $h(P)$ is a point, $P$ is a projective fiber subspace, and $S$ is a
smooth projective surface.  Take a surface minimal model $\rho:S\to S_0$.
If $K_{S_0}$ is nef, use general ample curves on $S_0$.  Otherwise the
classification of smooth projective surfaces gives ruling fibers with
canonical degree $-2$, or lines on $\mathbb P^2$ with canonical degree
$-3$.  Their general strict transforms $F$ satisfy $K_S\cdot F\geq-3$,
because $K_S-\rho^*K_{S_0}$ is effective.  In all cases the curves cover
$S$, so $F$ may be chosen not contained in the finitely many excluded
divisors.  This is the surface argument also used in
\cite[Lemma~3.3]{CT2023}.

The curve $F$ maps birationally to a curve $C\subset P$: it is not
contained in the exceptional or conductor locus.  By the projection
formula and \eqref{eq:surface-effective-correction},
\[
 (K_U+P)\cdot C=(K_S+\Delta_S)\cdot F
                  \geq K_S\cdot F\geq-3.
\]
Its choice also ensures $C\not\subset\Supp D_j$ for every $j$.
\end{proof}

The surface estimate bounds the denominator introduced by extracting one
exceptional place of log discrepancy at most one.  We now turn it into a
bound for every global rank-one reflexive sheaf.

\begin{lemma}\label{lem:index-bound}
For every integer $n\geq0$ and $\varepsilon>0$, there is an integer
$I(n,\varepsilon)>0$ with the following property.  Let
$(T,B+\bM)$ be a rational generalized klt compact K\"ahler threefold pair,
with $T$ globally strongly $\Q$-factorial.  Suppose that
\begin{enumerate}[label=\textup{(\roman*)}]
 \item every log discrepancy is at least $\varepsilon$;
 \item at most $n$ exceptional valuations have log discrepancy at most one;
 \item every exceptional log discrepancy greater than one is at least
       $1+\varepsilon$.
\end{enumerate}
Then $\cF^{[I(n,\varepsilon)]}$ is invertible for every global coherent
rank-one reflexive sheaf $\cF$ on $T$.
\end{lemma}

\begin{proof}
We induct on the number of exceptional places with log discrepancy at
most one.  The terminal case is controlled by local canonical indices.
For the induction step, extract one place, bound its intersection degree
using the preceding surface estimate, and descend an integral line bundle.

Replace $\varepsilon$ by $\min\{\varepsilon,1\}$, so assume
$0<\varepsilon\leq1$.  We construct integers divisible by those from the
preceding induction stage.  When $n=0$, the pair is generalized terminal.
The underlying space is ordinary terminal and
\[
 a(E;T,0)\geq a(E;T,B+\bM)\geq1+\varepsilon
\]
for every exceptional valuation.  By
Lemma~\ref{lem:terminal-local-index} and the observation following that
lemma, every canonical point index $r>1$ satisfies
$r\leq1/\varepsilon$.  Each local representative of $\cF$ is
$\Q$-Cartier by the global strong condition, and its Cartier index divides
$r$.  Thus one may take
\[
 I(0,\varepsilon)=\operatorname{lcm}
                     \{1,\ldots,\lfloor1/\varepsilon\rfloor\}.
\]
Local invertibility of this reflexive power at every point is precisely
global invertibility; no trivialization on all of $T$ is asserted.

Suppose $n>0$ and put $I'=I(n-1,\varepsilon)$.  If there is no exceptional
valuation of discrepancy at most one, the preceding case applies.  Otherwise
choose one such place, of discrepancy $a\in[\varepsilon,1]$, and apply
Proposition~\ref{prop:extraction} to obtain a projective extraction
$h:U\to T$ with that single exceptional prime $P$.  Write its effective
crepant boundary as
\[
 B_U=B_U^{\mathrm{str}}+(1-a)P.
\]
Crepancy preserves all discrepancies and lowers the number of exceptional
low places by one.  The induction hypotheses therefore hold on $U$, so
$I'$ clears the index of every global rank-one reflexive sheaf on $U$.

First, $-P$ is $h$-ample.  If $H$ is an $h$-ample line bundle, its reflexive
pushforward is rational Cartier on $T$.  Lemma~\ref{lem:trace} gives
\[
 H=h^*H_T+cP
\]
as a rational bundle comparison.  The divisor $cP$ is $h$-ample and
$h$-exceptional, hence is nonpositive by negativity.  Since $h$ has a
positive-dimensional fiber, $c\neq0$, and therefore $c<0$.

The same trace comparison for the nef data gives
\[
 M_U=h^*M_T-sP,\qquad s\geq0.
\]
Apply Lemma~\ref{lem:exceptional-surface-curve} with the components of
$B_U^{\mathrm{str}}$ as the excluded divisors.  For its contracted curve
$C\subset P$, the other boundary has nonnegative degree and
$M_U\cdot C=-sP\cdot C\geq0$.  Crepancy yields
\[
 -3\leq(K_U+P+B_U^{\mathrm{str}}+M_U)\cdot C
         =aP\cdot C<0.
\]
Consequently
\begin{equation}\label{eq:exceptional-degree-bound}
 0<-P\cdot C\leq\frac{3}{\varepsilon}.
\end{equation}

Now let $\cF$ be any global rank-one reflexive sheaf on $T$, and let
$\cF_U$ be its reflexive strict transform.  In additive rational-bundle
notation, Lemma~\ref{lem:trace} gives
\begin{equation}\label{eq:sheaf-extraction-comparison}
 h^*\cF=\cF_U+qP,\qquad q\in\Q.
\end{equation}
This comparison is defined by local meromorphic frames agreeing off the
exceptional locus, and requires no global meromorphic section of $\cF$.
The sheaves $\cF_U^{[I']}$ and $\cO_U(I'P)$ are line bundles.  Intersecting
\eqref{eq:sheaf-extraction-comparison} with $C$ gives
\[
 q=\frac{I'\cF_U\cdot C}{-I'P\cdot C},
 \qquad k:=-I'P\cdot C\in\Z,
 \quad 1\leq k\leq\frac{3I'}{\varepsilon}.
\]
The numerator is an integer.  Define
\begin{equation}\label{eq:index-recurrence}
 I(n,\varepsilon)
   =I'\operatorname{lcm}
       \{1,\ldots,\lceil3I'/\varepsilon\rceil\}.
\end{equation}
Then $I:=I(n,\varepsilon)$ is divisible by $I'$ and $Iq\in I'\Z$.
Thus $I\cF_U+(Iq)P$ represents an integral line bundle $L$ on $U$,
with its prescribed comparison to $h^*\cF^{[I]}$ off the exceptional
locus.  Equation~\eqref{eq:sheaf-extraction-comparison} shows that $L$ has
degree zero on every $h$-contracted curve.

Increase the coefficient $1-a$ of $P$ by a sufficiently small positive
rational number $\delta$.  On a fixed log resolution carrying the data,
the inequalities defining generalized klt remain strict for all sufficiently
small $\delta$, so the new pair is effective and generalized klt.  Its
adjoint is
\[
 h^*(K_T+B+M_T)+\delta P,
\]
which is $h$-antiample.  Lemma~\ref{lem:descent} now descends $L$
\emph{integrally}.  More explicitly, local ordinary klt replacement and
the analytic base-point-free theorem generate $L^m$ for every sufficiently
large integer $m$; descend two consecutive powers and take their quotient.
No further index is introduced.  The descended line bundle agrees with
$\cF^{[I]}$ away from $h(P)$, a set of codimension at least two in $T$.
Reflexive extension identifies them on $T$.  This proves the induction,
including the case $a=1$.
\end{proof}

\subsection{Nef endpoints and terminal lifting}

The uniform index bound will make the low discrepancies constant on a
tail of any small sequence.  After extracting those places, we must compare
each lifted relative program with the prescribed positive model downstairs.
The next Lemma provides both the discrepancy comparison during the run
and the crepant morphism at its nef endpoint.  It will also be used for
strata and threshold lifts.

\begin{lemma}\label{lem:nef-end-comparison}
Let $Y_0,Y,T_+$ be normal compact spaces projective and bimeromorphic over
a normal space $S$, and let $W$ be a common smooth modification
projective over each of the three spaces, with maps
$p_0:W\to Y_0$, $p:W\to Y$, and $q:W\to T_+$.
Let $A_0,A,A_+$ be rational line bundles on the respective spaces.
Suppose their prescribed meromorphic comparisons give
\begin{equation}\label{eq:two-target-comparison}
 P_0=P+G=P_++F,
 \qquad P_0=p_0^*A_0,\quad P=p^*A,\quad P_+=q^*A_+,
\end{equation}
where $G$ is effective and $p$-exceptional, and $F$ is effective and
$q$-exceptional.  If $A_+$ is nef over $S$, then $F-G\geq0$.
If also $A$ is nef over $S$, then $P=P_+$.  If, in addition, $A_+$ is
ample over $S$, the induced bimeromorphic map $Y\dashrightarrow T_+$
is a projective morphism $g$, and $A=g^*A_+$ as rational line bundles
with the given comparison.
\end{lemma}

\begin{proof}
Set $E=P-P_+=F-G$.  Over $Y$, the divisor $-E=P_+-P$ is nef, while
\[
 p_*(-E)=-p_*F\leq0,
\]
since $G$ is $p$-exceptional.  Lemma~\ref{lem:negativity} gives
$-E\leq0$, proving the first assertion.  If $A$ is nef over $S$, then
$E$ is nef over $T_+$ and
\[
 q_*E=-q_*G\leq0,
\]
since $F$ is $q$-exceptional.  Negativity gives $E\leq0$, and hence
$E=0$.  The two exceptional supports need not coincide.

Assume $A_+$ is relatively ample.  A fiber of $p$ maps by $q$ into a
projective fiber of $T_+\to S$.  Equality $p^*A=q^*A_+$ gives degree zero
on every curve in that fiber.  A nonconstant image would be detected by a
curve of positive degree against $A_+$, so $q$ is constant on each
connected fiber of $p$.  Proper factorization through the normal space
$Y$ gives $q=g\circ p$.  The map $g$ is projective because $Y\to S$ is
projective and $T_+\to S$ is separated.  Finally, pullback by $p$ is
injective on rational line bundles, by $p_*\cO_W=\cO_Y$ and the
projection formula; therefore $A=g^*A_+$.  This proves the actual bundle
identity, not merely numerical equality.  Compare
\cite[Lemma~2.12]{LMT2022} for the algebraic form of this comparison.
\end{proof}

\begin{theorem}\label{thm:gklt-three}
Let $(X_0,B_0+\bM)$ be a rational generalized klt pair of dimension at
most three on a normal globally strongly $\Q$-factorial compact K\"ahler
space, with fixed rational b-line data as in
Definition~\ref{def:generalized-data}.  Every sequence of elementary
birational steps in the sense of
Section~\ref{sec:relative} starting from this pair terminates.  The
compact K\"ahler contraction bases may vary.
\end{theorem}

\begin{proof}
In dimension at most two there are no nontrivial small birational
contractions, and Lemma~\ref{lem:cycle-loss} gives the result.  In dimension
three that lemma first removes a finite divisorial prefix.  Relabel the
remaining small sequence as
\[
 (X_i,B_i+\bM)\dashrightarrow(X_{i+1},B_{i+1}+\bM),\qquad i\geq1.
\]
All boundaries are now strict transforms of $B_1$.

By Lemma~\ref{lem:low-places}, there is $\varepsilon_0>0$ such that the
initial set of exceptional valuations with discrepancy below
$1+\varepsilon_0$ is finite.  Since discrepancies increase and small maps
preserve exceptionality, the sets
\[
 \mathcal L_i=\{E\text{ exceptional over }X_i:
                       a(E;X_i,B_i+\bM)\leq1\}
\]
form a decreasing sequence of finite sets.  After discarding a finite
prefix, they equal one set $\mathcal L=\{E_1,\ldots,E_n\}$.
At the first model of this tail, the finitely many exceptional values in
$(1,1+\varepsilon_0)$ have a positive minimum distance from one; use
$\varepsilon_0$ itself if there are no such values.  Choose
$\varepsilon>0$ no larger than this distance, than $\varepsilon_0$, than
one, and than the positive lower bound for all discrepancies of the initial
gklt pair furnished by Lemma~\ref{lem:low-places}.  Monotonicity and the
stabilization of $\mathcal L$ show that, on every model of this tail,
all discrepancies are at least $\varepsilon$, and every exceptional
discrepancy greater than one is at least $1+\varepsilon$.

Apply Lemma~\ref{lem:index-bound} with these fixed $n$ and
$\varepsilon$.  Let $r$ clear the coefficients of $B_1$ and let $p$ clear
the higher nef datum.  For example
\[
 J=I(n,\varepsilon)\operatorname{lcm}(r,p)
\]
makes each $J(K_{X_i}+B_i+M_{X_i})$ an integral line bundle, and makes
$J\bM$ integral on every higher carrier.  Comparing the actual adjoint
bundles on a smooth carrier shows that every discrepancy belongs to
$J^{-1}\Z$: the canonical bundle there is integral and the coefficients
of the corresponding integral divisor difference are integers.  The same
argument applies after any further resolution.  Thus each nondecreasing
sequence
\[
 a(E_\ell;X_i,B_i+\bM)\in J^{-1}\Z\cap[\varepsilon,1]
\]
is eventually constant.  After another finite prefix, all these values
are constant simultaneously.

Use Proposition~\ref{prop:extraction} to extract exactly
$E_1,\ldots,E_n$ over the first model of this final tail:
\[
 h_1:(Y_1,\Delta_1+\bM)\longrightarrow(X_1,B_1+\bM).
\]
The boundary is effective and crepant.  Every valuation exceptional over
$Y_1$ is exceptional over $X_1$ and is not in $\mathcal L$, so its
discrepancy is greater than one.  Hence this is a globally strongly
$\Q$-factorial generalized terminal model.  If $n=0$, take $h_1=\id$.

Consider the first remaining flip over $Z_1$.  Starting from $Y_1$, use
Proposition~\ref{prop:relative-steps} to take relative steps over $Z_1$
while the adjoint is not nef.  We show that no divisorial contraction can
occur in this run.  At any finite stage $Y_j$, a common smooth model gives
\[
 P_0=P_j+G_j=P_++F,
\]
where $P_0$ is pulled back from $Y_1$, $P_j$ from the current lifted
adjoint, and $P_+$ from $K_{X_2}+B_2+M_{X_2}$.  These are the prescribed
adjoint comparisons.  The accumulated error $G_j$ is effective and
exceptional over $Y_j$; the error $F$ of the original flip is effective
and exceptional over $X_2$.  The positive adjoint is ample over $Z_1$.
The first assertion of Lemma~\ref{lem:nef-end-comparison}, which does not
require nefness on $Y_j$, gives $G_j\leq F$.

If a divisor were lost, its strict transform from $Y_1$ would have a
positive coefficient in $G_j$, by the strict discrepancy comparison at
its contraction.  But every divisor on $Y_1$ is either the transform of a
divisor on $X_1$, or one of the $E_\ell$.  In the first case its
coefficient in $F$ is zero because the original map is small.  In the
second case that coefficient is
\[
 a(E_\ell;X_2,B_2+\bM)-a(E_\ell;X_1,B_1+\bM)=0
\]
by stabilization.  Both cases contradict $G_j\leq F$.  Thus the entire
lifted run is small.  It stays generalized terminal, and
Proposition~\ref{prop:terminal-three} forces it to reach a nef endpoint
after finitely many steps.

At that endpoint, Lemma~\ref{lem:nef-end-comparison} gives equality of
the two adjoint pullbacks and a projective crepant morphism
\[
 h_2:(Y_2,\Delta_2+\bM)\longrightarrow(X_2,B_2+\bM).
\]
Here $Y_2$ denotes the endpoint, rather than the first intermediate model
of the finite run.  Since the lifted transformations were small, the
exceptional primes of $h_2$ are exactly the same $E_\ell$, with the same
boundary coefficients.  We may therefore repeat this construction over
each successive base $Z_i$.

Each nontrivial flip below forces at least one step above.  Indeed, take a
negative curve $C\subset X_i$ over $Z_i$.  Projectivity of $h_i$ supplies
a curve $\Gamma\subset Y_i$ mapping onto $C$ with positive degree.  The
crepant identity gives
\[
 (K_{Y_i}+\Delta_i+M_{Y_i})\cdot\Gamma
   =\deg(\Gamma/C)(K_{X_i}+B_i+M_{X_i})\cdot C<0.
\]
Thus the starting lifted adjoint is not nef over $Z_i$.
An infinite original sequence would consequently concatenate to an
infinite sequence of small generalized terminal elementary steps, with
fixed rational boundary coefficients and the same rational nef b-line
data.  All spaces remain compact K\"ahler and globally strongly
$\Q$-factorial.  The bases may change, exactly as allowed in
Proposition~\ref{prop:terminal-three}, which gives the contradiction.
\end{proof}
